\section{Limitations and threats to validity}\label{sec:limitations}

% The one place where the caveats of the paper are stated. Other sections refer here.

% The finite-element reference (Section 5.4, sections/fem-reference.tex) covers disk and bearing
% only; widen this paragraph if another external reference is added.
\paragraph{One finite-element reference, no experimental one.} Every reference of the battery
is analytic. Class A references are exact for rigid bodies; class B references assume a
pressure distribution; class C references are the foundation's own. For two class B scenarios,
\scn{disk} and \scn{bearing}, Section~\ref{sec:fem} places a linear-elastic continuum in the
band. It settles where a continuum support falls there (on the Winkler cosine for the thin
sleeve, at the rigid-punch end for the battery's thick slab) and that on the thick slab the
methods that integrate the foundation correctly are off by the distance between a Winkler
foundation and a continuum, which no discretization removes. It does not settle more. It is a
static, small-strain solution of the frictionless normal problem, exact for the spinning punch
and the $\mu\to0$ limit for the sleeve, and it assumes sliding everywhere, while the runs are
dynamic, damped and pass through stick; the support's Poisson ratio is not given by the
battery, and $E$ follows from $K$ through the thin-layer mapping. It covers two scenarios:
\scn{ring}, \scn{contensou} and \scn{cone} were not computed, and no class A or C scenario,
impact, rolling resistance or flutter has a continuum reference. Apart from the thin layers of
Section~\ref{sec:fem}, where a continuum behaves as one, nothing in this paper shows that a
Winkler foundation of stiffness $K$ describes a physical contact, and no measurement was used. Conformal cylindrical contact has finite-element and analytic references in the
literature on clearance joints \cite{Liu2007,Pereira2011,Pereira2015}, and patch friction has
experimental ones \cite{Kudra2015}; we cite those works for what their titles and abstracts
state and used none of their reference solutions.

\paragraph{A regime chosen for the laws.} The stiffness of every scenario was chosen so that
the explicit coupling is stable at the base step of 1\,ms, with two exceptions written to show
the limit (\scn{tower}, \scn{massratio}). The foundations are therefore soft, $10^6$ to
$10^8$\,N/m$^3$, orders of magnitude below engineering materials over a layer of millimetres,
where the laws would need a step proportional to $K^{-1/2}$ and solvers implicit in the contact
would not; one decade of stiffness was explored, at one step, and the rigid limit of class A
was not tested. The damping factor $h = 5$\,s/m has no physical justification and was not tied
to a coefficient of restitution; with $K$ it sets the errors of \scn{bowl}, \scn{trough},
\scn{turntable} and \scn{edge}. The stiction constants $d$ and $\tau$ were not swept. The
battery did not record in which runs the step bounds of Section~\ref{sec:coupling} were active;
from the loads, they are in \scn{tower}, \scn{massratio} and \scn{stack} at the base step and
not where one body rests under its own weight. Where they act, what is benchmarked is the law
regularized by the step, and a flat step series is evidence about the law as given only where
they do not.

\paragraph{Few, hand-picked scenarios that the laws have seen.} The
\result{battery/scenarios} scenarios are not a sample of the problems that a user will meet;
several test the same mechanism, and class B has \result{battery/n/B}. The laws were developed
on \result{battery/dev} of them and changed in response; the other \result{battery/heldout},
the second set, were written before any law was run on them, but three changes to the laws were
adopted after their first run (Section~S4.3), so no part of the battery is an out-of-sample
test of the two laws. Every change was general (no scenario has a constant of its own). The
ratios, intervals and $p$-values of Section~\ref{sec:ratios} are therefore descriptive: the
bootstrap treats as exchangeable units that are not, and the family of primary comparisons and
the rules for capped and failed runs were fixed after the data had been seen (Section~S5). The
battery records the class B errors without sign (Table~\ref{tab:fem} signs those of \scn{disk}
and \scn{bearing} from their stop times). The findings that we expect to survive another
battery are those tied to a mechanism, not the counts.

\paragraph{The native configurations and the baselines are ours.} No configuration of a native
model was reviewed by the developers of its engine; they were built from the documentation,
revised after first results on the development scenarios, and contain asymmetries (Section~S3).
The main text therefore claims of the native models only what follows from their model family.
Their other results, several of them unexplained and several scattered over nearby starts, are
in Section~S3 as results of our configurations and should not be read as the performance of
those engines in expert hands. The vertex rule follows our reading of Hasegawa and Sato
\cite{Hasegawa2004}; the paper was not accessible to us again, so the rule was not re-checked
against its equations. \law{pcm} departs from its publications (Section~S1.13), and the
published models on their own meshes were not run. The statements of
Section~\ref{sec:related} about the models of Hasegawa and Fujii, Luo and Nahon, and Wakisaka
and Sugihara rest on the classification of Elandt et al.~\cite{Elandt2019}, not on the primary
sources.

\paragraph{Limits of validity of the laws.} (i) \emph{Explicit coupling}: the normal stiffness
is bounded by Eq.~\eqref{eq:stability}, and above a load threshold the step bounds apply a
softer stick, a wider stick--slip transition and less damping than the parameters say; the
remedy is an implicit coupling with the Jacobians of the wrench. (ii) \emph{Friction-induced
flutter}: beyond Eq.~\eqref{eq:flutter} a sliding face under a body with a raised centre of
mass is unstable in rocking; the ratio of the two sides grows as $\mu\,z_c\sqrt{K}/h$, the
class B scenarios are below the limit at their $K$ and above it at $10K$, and a damping stated
as a time constant, which would reverse that dependence, would change the published law and
every result and was not tried. (iii) \emph{One normal per pair of convex pieces}
\cite{Elandt2019}: non-convex bodies must be decomposed exactly, curved conforming contacts are
represented facet by facet, and \scn{journal} does not converge. (iv) \emph{Energy}: Gonthier's
normal force is not conservative on faceted or decomposed geometry, and the elastic force of
\law{distributed} is not where the pair normal turns with the overlap. (v) \emph{What ``exact''
covers}: the geometry and the normal wrench, for the polyhedra given and not for the curved
bodies that they approximate; the friction is exact for a linear pressure on each flat triangle
and approximate elsewhere, and the stiction is a lumped regularization with two constants that
are not material properties (Sections~S1.8, S1.9 and~S1.11).

\paragraph{Unexplained observations.} \law{pcm}, whose normal law, damping and friction lever
are nominally those of the two laws, did not flutter on the disk up to $100K$. The value of the
energy measure of \law{gonthier} at the base step differs from that of a single run made when
the scenario was written, and the change that moved it was not identified. How much of the
growth of the errors of \scn{cone} and \scn{bearing} at $10K$ is flutter and how much the sag
of the facets was not separated. Section~S5.5 lists the runs that would settle these and the
other open points.

\paragraph{Timings.} The costs are those of one processor, one thread configuration and one
step size; the two friction quadratures were not timed, Chrono SMC is sub-stepped by its
backend, and the native models were run at the step of the battery, not at the larger steps
that implicit solvers allow.
